\documentclass[12pt,a4paper,oneside]{report}
 
\usepackage[utf8]{inputenc}
\usepackage[T1]{fontenc}
\usepackage[italian,english]{babel}
\usepackage{geometry}
\usepackage{setspace}
\usepackage{titlesec}
\usepackage{hyperref}
 
\geometry{
    top=3cm,
    bottom=3cm,
    left=3cm,
    right=3cm
}
 
\setstretch{1.2}
 
\title{\textbf{Quantum Annealing Learning Search per la risoluzione di problemi QUBO} \\[0.3cm]
\large Anteprima della struttura della tesi}
\date{}
 
\begin{document}
 
\maketitle
 
\tableofcontents
 
\chapter{Introduzione}
 
\section{Motivazioni}
Spiegare perché studiare QALS e il quantum annealing applicato a problemi QUBO;
contesto dell'informatica quantistica e limiti dei metodi classici che motivano
l'approccio ibrido.
 
\section{Obiettivi}
Definire cosa si vuole dimostrare/verificare nella tesi: analizzare le prestazioni
di QALS sulla variante 0/1 del Knapsack Problem e proporre miglioramenti.
 
\chapter{Background teorico}
 
\section{Quantum Annealing}
Descrivere il principio fisico del quantum annealing, l'hardware D-Wave e 
le topologie (Chimera, Pegasus) con i relativi limiti di embedding.
 
\chapter{Quantum Annealing Learning Search (QALS)}
 
\section{Algoritmo}
Descrivere l'idea generale e lo pseudocodice di QALS (tabù matrix, parametri.)
 
\chapter{Knapsack Problem}
 
\section{Formulazione QUBO}
Definire le variante del Knapsack Problem trattata (0/1) e introdurre il 
termine di penalità quadratico per trasformarle in problemi QUBO.
 
\section{Derivazione della matrice Q}
Sviluppare passo per passo il calcolo algebrico che porta dalla funzione
obiettivo penalizzata alla matrice $Q$, distinguendo i termini lineari e
quadratici.
 
\section{Scelte del parametro \(\lambda\)}
Confrontare i tre metodi di calcolo di \(\lambda\) utilizzati nei report
(\(\sum p_i / C\), \(\sum p_i / 3\), \(\sum p_i / (650 \cdot C)\)) e il loro effetto
sulla penalizzazione del vincolo.
 
\chapter{Metodologia sperimentale}
 
\section{Scaling QUBO}
Spiegare l'effetto della moltiplicazione della matrice \(Q\) per un fattore
\(\alpha\).
 
\section{Parametro Swap(k, TIMES)}
Intro distinzione tra il parametro \(k\) (numero di chiamate
all'annealer) e TIMES (numero di ripetizioni di QALS/annealer), motivandone
l'introduzione e i risultati preliminari.
 
\section{Confronto con Gurobi}
Presentare la validazione delle soluzioni trovate da QALS confrontandole con
quelle ottenute da Gurobi sulle istanze di test.
 
\chapter{Diagnostica di QALS}
 
\section{Distanze di Hamming}
Definire le quattro metriche di distanza di Hamming usate per analizzare il
comportamento di QALS.
 
\section{Clustering agglomerativo}
Descrivere la costruzione della matrice delle distanze, il clustering
agglomerativo delle soluzioni esplorate, l'individuazione del medoid e
l'interpretazione in termini di concentrazione/dispersione della ricerca.
 
\section{Decay factor}
Presentare l'introduzione del decay factor nella matrice tabù, la sua
formulazione matematica e gli effetti osservati sui risultati sperimentali.
 
\chapter{Proposta di miglioramento}
Illustrare le criticità individuate nella diagnostica, la modifica proposta 
all'algoritmo, la sua motivazione teorica e il confronto con la versione originale.
 
\chapter{Risultati sperimentali}
 
\section{Riepilogo comparativo per istanza}
Riassumere e confrontare i risultati ottenuti su tutte le istanze di test.
 
\chapter{Conclusioni}
Riassumere il lavoro svolto, i risultati principali, i limiti, i possibili 
sviluppi futuri e applicazioni.
 
 
\end{document}
